\documentclass[aspectratio=169,10pt]{beamer}
\usetheme[numbering=fraction,progressbar=frametitle]{metropolis}
\usepackage{fontspec}
\setsansfont{Latin Modern Sans}
\usepackage{amsmath,amssymb,mathtools,graphicx,booktabs}
\usepackage{polyglossia}
\setdefaultlanguage{english}
\setbeamertemplate{navigation symbols}{}
\setbeamerfont{frametitle}{size=\large}
\setbeamerfont{normal text}{size=\normalsize}
\setlength{\parskip}{.45em}
\hypersetup{pdfauthor={},pdfsubject={Lecture slides},pdfkeywords={},pdftitle={Orthogonality and distances}}
\providecommand{\N}{\mathbb{N}}
\providecommand{\Z}{\mathbb{Z}}
\providecommand{\D}{\mathbb{D}}
\providecommand{\Q}{\mathbb{Q}}
\providecommand{\R}{\mathbb{R}}
\providecommand{\C}{\mathbb{C}}
\providecommand{\vect}[1]{\overrightarrow{#1}}
\providecommand{\vi}{\vec{\imath}}
\providecommand{\vj}{\vec{\jmath}}
\providecommand{\vk}{\vec{k}}
\providecommand{\norm}[1]{\left\lVert #1\right\rVert}
\providecommand{\vc}[2]{\begin{pmatrix}#1\\#2\end{pmatrix}}
\providecommand{\deter}[2]{\det\left(#1,#2\right)}
\providecommand{\abs}[1]{\left\lvert #1\right\rvert}
\providecommand{\dd}{\,\mathrm{d}}
\providecommand{\e}{\mathrm{e}}
\providecommand{\Card}{\operatorname{Card}}
\providecommand{\E}{\operatorname{E}}
\providecommand{\V}{\operatorname{V}}
\providecommand{\degre}{^{\circ}}
\providecommand{\un}[1]{\,\mathrm{#1}}
\providecommand{\sci}[2]{#1\times10^{#2}}
\providecommand{\pgcd}{\operatorname{PGCD}}
\providecommand{\ppcm}{\operatorname{PPCM}}
\title{Orthogonality and distances}
\subtitle{Lecture and worked examples}
\date{}
\author{}
\begin{document}
\begin{frame}[plain]
\vfill
{\usebeamerfont{title}\inserttitle\par}
\medskip
{\usebeamerfont{subtitle}\insertsubtitle\par}
\medskip
{\usebeamercolor[fg]{progress bar}\rule{\textwidth}{.5pt}}
\vfill
\end{frame}
\begin{frame}{Objectives}
Calculate a dot product; justify an orthogonality; find an orthogonal projection and a distance; calculate a volume.
\end{frame}
\begin{frame}{Course notes}
The coordinate calculations in this chapter are carried out in an orthonormal frame of space.
\end{frame}
\begin{frame}{Dot product in space}
For non-zero vectors \(\vec u,\vec v\), \(\vec u\cdot\vec v=\|\vec u\|\|\vec v\|\cos\theta\), where \(\theta\in[0,\pi]\) is their angle. If one of them is zero, the dot product is \(0\).
\[\vec u\cdot\vec v=\tfrac12(\|\vec u+\vec v\|^2-\|\vec u\|^2-\|\vec v\|^2).
\]
In an orthonormal frame, for \(\vec u=(x,y,z)\) and \(\vec v=(x',y',z')\),
\[\vec u\cdot\vec v=xx'+yy'+zz'.\]
The dot product is symmetric and bilinear; \(\vec u\cdot\vec u=\|\vec u\|^2\). Two vectors are orthogonal if and only if their dot product is zero.
\end{frame}
\begin{frame}{Dot product in space}
\(\vec u=(1,2,2)\) and \(\vec v=(2,-1,0)\) are orthogonal: \(2-2+0=0\). Their norms are \(3\) and \(\sqrt5\).
\end{frame}
\begin{frame}{Line, plane and normal vector}
Two lines are orthogonal if their direction vectors are. A line is orthogonal to a plane if it is orthogonal to two non-collinear directions of that plane. A normal vector to a plane is a non-zero vector orthogonal to all its directions.
\end{frame}
\begin{frame}{Line, plane and normal vector}
If \(\vec n\cdot\vec u=\vec n\cdot\vec v=0\), then, for every \(\vec w=s\vec u+t\vec v\), \(\vec n\cdot\vec w=s\vec n\cdot\vec u+t\vec n\cdot\vec v=0\).
\end{frame}
\begin{frame}{Line, plane and normal vector}
If \(A\) belongs to the plane \(\mathcal P\) with normal vector \(\vec n\), then \(M\in\mathcal P\) if and only if \(\vect{AM}\cdot\vec n=0\).
\end{frame}
\begin{frame}{Orthogonal projection and distance}
The orthogonal projection \(H\) of \(M\) onto a plane is the unique point of the plane such that \(\vect{MH}\) is normal to the plane; if \(M\) belongs to the plane, \(H=M\). The distance from \(M\) to the plane is \(MH\). For a line with direction vector \(\vec u\), the projection satisfies \(H\in d\) and \(\vect{MH}\cdot\vec u=0\).
\end{frame}
\begin{frame}{Orthogonal projection and distance}
For every point \(P\) of the plane, \(\vect{MH}\perp\vect{HP}\), so \(MP^2=MH^2+HP^2\ge MH^2\). Equality holds only for \(P=H\). The projection is the closest point of the plane.
\end{frame}
\begin{frame}{Orthogonal projection and distance}
If \(d=A+\R\vec u\), write \(H=A+t\vec u\) and solve \((A-M+t\vec u)\cdot\vec u=0\). Hence \(t=\dfrac{\vect{AM}\cdot\vec{u}}{\lVert\vec u\rVert^2}\).
\end{frame}
\begin{frame}{Angles, areas and volumes}
For an angle \(\widehat{ABC}\), use \(\vect{BA}\) and \(\vect{BC}\): \(\cos \widehat{ABC}=\dfrac{\vect{BA}\cdot\vect{BC}}{BA\times BC}\). For the acute angle between two lines, take the absolute value of the quotient.
The area of a triangle is \(\dfrac{bh}{2}\); the volume of a tetrahedron is \(\dfrac{\mathcal{A}_{\rm base}h}{3}\), where \(h\) is an orthogonal distance. The units are squared and cubed respectively.
\end{frame}
\begin{frame}{Expansions and loci}
By bilinearity,
\[\|\vec u+\vec v\|^2=\|\vec u\|^2+2\vec u\cdot\vec v+\|\vec v\|^2,\qquad
\|\vec u-\vec v\|^2=\|\vec u\|^2-2\vec u\cdot\vec v+\|\vec v\|^2.\]
So we can also calculate \(\vec u\cdot\vec v=\dfrac12(\|\vec u\|^2+\|\vec v\|^2-\|\vec u-\vec v\|^2)\).
For two distinct points \(A,B\) with midpoint \(I\), the locus of points \(M\) such that \(MA=MB\) is the plane through \(I\) with normal vector \(\vect{AB}\): the \emph{perpendicular bisector plane}. Indeed, writing the vectors from \(I\), the expansions give \(MA^2-MB^2=2\vect{IM}\cdot\vect{AB}\). Equality of the distances is therefore equivalent to \(\vect{IM}\cdot\vect{AB}=0\).
\textbf{Extension, not specifically assessed in 2026-2027.} Two planes are perpendicular when their normal vectors are orthogonal. This property must not be confused with parallelism of planes, which corresponds to collinear normal vectors.
\end{frame}
\begin{frame}{Research activity}
\textbf{The closest point on a line - 35 min.}

In an orthonormal frame of space, consider \(A(1;0;0)\), \(M(0;2;3)\) and the line \(d\) through \(A\) with direction vector \(\vec u=(1;1;0)\). We look for the point of \(d\) closest to \(M\).
\end{frame}
\begin{frame}{Research activity (2)}
\begin{enumerate}
\setcounter{enumi}{0}\item Express, in terms of a real parameter \(t\), the coordinates of a point \(H\) of \(d\).

\item Find the value of \(t\) for which \(\vect{MH}\cdot\vec u=0\). Deduce the coordinates of the resulting point \(H\).

\item Calculate the distance \(MH\).

\item Let \(P\) be any point of \(d\). Prove that \(MP^2=MH^2+HP^2\), then deduce that \(H\) is the only point of \(d\) closest to \(M\).

\end{enumerate}
\end{frame}
\end{document}
